% problem-000024 3 金属氧化物分离与分析
\section*{3. 金属氧化物分离与分析}
某⾦属M的氧化物A是易挥发的液体，有毒，微溶于⽔。在氯化氢⽓氛中蒸发A的盐酸溶液（含1.000 g的 A）得到 1.584 g 晶体 $\mathbf { M C l } _ { 3 } { \cdot } 3 \mathrm { H } _ { 2 } \mathrm { O } .$ o

\subsection*{3-1}
假定转化反应按化学计量⽐进⾏，通过计算确定A的化学式，并写出⾦属M的价电⼦构型。

从精炼镍的阳极泥提取M的步骤为：

\subsection*{(1)}
⽤王⽔处理溶解 Pt、Pd 和 Au；

\subsection*{(2)}
固体不溶物与碳酸铅共热，然后⽤稀硝酸处理除去可溶物B；

\subsection*{(3)}
剩余固体不溶物与硫酸氢钠共熔，⽤⽔浸取除去可溶物 $\mathrm { c } ;$

\subsection*{(4)}
留下的固体不溶物与过氧化钠共熔，⽤⽔浸取，过滤。滤液中含有M的盐D，D为 $\mathrm { N a _ { 2 } S O _ { 4 } }$ 型电解质，⽆⽔D 中氧含量为30.32%；

\subsection*{(5)}
往D溶液中通⼊Cl_{2}，加热，收集 $4 0 – 5 0 ^ { \circ } \mathrm { C }$ 的馏分，得到A。A收集在盐酸中，加热得到E的溶液，E中Cl 含量为 67.14%；

$
\mathrm{NH} _ {4} \mathrm{Cl}
$

$
\mathbf {F},
$

$
\mathbf {M} _ {\circ}
$

已 $\mathrm { I | \varphi _ { A } ^ { \ominus } ( N O _ { 3 } ^ { - } / N O ) = 0 . 9 5 7 \vee , ~ \varphi ^ { \ominus } ( [ P t C l _ { 4 } ] ^ { 2 - } / P t ) = 0 . 7 5 5 \vee , ~ \varphi ^ { \ominus } ( [ P t C l _ { 6 } ] ^ { 2 - } / [ P t C l _ { 4 } ] ^ { 2 - } ) = 0 . 6 8 0 ~ V , }$ ，所有离⼦活度系数为 1.00。

\subsection*{3-2}
写出步骤(1)中除去Pt的化学反应⽅程式，计算该反应在298 K下的 $K ^ { \ominus } { } _ { c }$

\subsection*{3-3}
写出步骤(2)中除去银的化学反应⽅程式。

\subsection*{3-4}
D、E、F中均只有1个⾦属原⼦，试通过计算确定D、E、F的化学式。
